\section{Operator-Based Optimization and Design Recovery}
\label{sec:solver_vm}\label{sec:ocs}\label{sec:vm_solver}\label{sec:mccormick_blowup}
We combine OCS working solves, structured primal improvements and complete-domain
bounds for the continuous LP, then search and replay integral designs.
These stages share coefficients but have distinct stopping criteria.

\paragraph{Operators replace matrices.}
For a working LP $\min_x f(x)+g(Kx)$, $f(x)=c^Tx+I_{[\ell,u]}(x)$
and $g$ encodes row equalities and inequalities. An operator computes $Kx$
without storing its matrix; its adjoint $K^Ty$ accumulates row prices back to
coordinates, including shared hardware and capacity. Generated incidence and
coefficient reuse preserve independent invocation decisions.

\paragraph{GPU first-order working solves.}
Our project copy of OCS 0.3.2 uses centered Halpern primal--dual hybrid gradient
updates with fixed diagonal steps and block-average restarts
\citep{chambolle2011first,halpern1967fixed,applegate2022practicallargescalelinearprogramming}.
Each step applies $K,K^T$, objective-shifted box projection and the appropriate
dual sign projection. JAX~\citep{jax2018github} compiles the update blocks.
Reversible power-of-two scaling reduces coefficient disparity; outward
row/column sums choose steps with $\|\Sigma^{1/2}KT^{1/2}\|_2\le0.9$.
All residual and bound checks return to original coordinates.

\paragraph{Working supports and implicit full-domain points.}
A checked lift embeds each restricted support into the source LP. One
expanded-residency support has $36{,}508$ columns and $313{,}061$ rows;
OCS reached a checked $0.0774\%$ working interval in $313.25$ seconds on an
A100, including preparation and checking. This alone is not a full-master bound.
For the full input, we improve independently identified owner retention,
original tiles and streamed-weight choices, then represent temporal fields
implicitly as weighted translations. GPU evaluation streams edge events while
the central affine envelope is summed analytically. Capacity is checked on the
resulting fractional point; no full event-expanded vector is allocated.

\paragraph{Complete-domain LP certificate.}
Coverage prices are checked against every implementation in the frozen generators.
Work-mass and service rows extend those inequalities to every original event;
the minimum of the 21 branch bounds gives the family lower bound.
Together with the admitted implicit primal, this yields
$L=734{,}095.710904$ ms and $U=734{,}720.389756$ ms, hence
$(U-L)/U=0.0850226646\%$: a certified $0.0850\%$ LP objective gap, rounded.
The certificate closes the continuous objective interval through specialized
primal transformations and algebraic bounds, not a claim that generic PDHG
iterates converged on the entire expansion
(Appendix~\ref{app:native_lp_certificate}).
This certificate concerns the frozen native instance; it is not transferred to
the strengthened physical branches in Eq.~\eqref{eq:master}.
A new source with the corrected BF16 chip instead gives a checked
$0.0846842\%$ interval over its complete implementation and temporal domain
(Appendix~\ref{app:corrected_native_lp}). Hardware is fixed in this result;
the numerical reconstruction engine has not converged over the joint master.

\paragraph{Hardware branch and bound.}
For integral recovery, a finite tree partitions array counts, vector lanes,
memory bandwidth, host channels and local-service units. Optimistic resources
and capacity-priced retention paths bound each region; budget-infeasible
regions or those unable to beat a replayed incumbent are pruned.
Promising leaves are evaluated with actual capacity and implementation costs.
The demonstrated tree covers a declared hardware grid and conditional software
choices, not every discrete selector of the master. It does not re-solve the
full temporal LP at every node. Appendix~\ref{app:native_solve_workflow}
describes the bounds and the unresolved schedule leaves.

\paragraph{Recovery and replay.}
Recovery selects one branch, an integral cover, hardware allocations and legal
storage/transfer actions. Analytical replay recomputes services at the selected
rates, dependency readiness and capacities, yielding the reported chip latency.
The LP certificate and recovered design are separate outputs: integer choices
do not remove continuous reciprocal/product relaxation error, and the
continuous LP interval is not a discrete-design optimality guarantee.
